\subsection{The structure of quasi-Galois extensions}

PROPOSITION 13. --- Let N be a quasi-Galois extension of K. We denote by \(N_{r}\) the field of invariants of the group of all K-automorphisms of N and by \(N_{s}\) the relative separable algebraic closure of K in N (V, p.~44). Then :

\begin{enumerate}
\def\labelenumi{\alph{enumi})}
\item
  \(\mathbf{N}_r\) is the relative \(p\) -radical closure of \(\mathbf{K}\) in \(\mathbf{N}\) (V, p.~25).
\item
  \(N_{s}\) is a Galois extension of K and every K-automorphism of \(N_{s}\) extends in a unique fashion to an \(N_{r}\) -automorphism of N.
\item
  The fields \(\mathbf{N}_r\) and \(\mathbf{N}_s\) are linearly disjoint over \(\mathbf{K}\) and we have \(\mathbf{N} = \mathbf{K}[\mathbf{N}_r \cup \mathbf{N}_s]\) ; in other words, the canonical homomorphism of \(\mathbf{N}_r \otimes_{\mathbf{K}} \mathbf{N}_s\) into \(\mathbf{N}\) is an isomorphism.
\end{enumerate}

Let \(\Omega\) be an algebraic closure of K containing N as subextension (V, p.~23, Th. 2). Every K-automorphism of \(\Omega\) induces an automorphism of N because N is quasi-Galois. Therefore every element of \(N_r\) is invariant under the group of K-automorphisms of \(\Omega\) , hence \(p\) -radical over K (V, p.~53, Cor. 3). Conversely every element of N which is \(p\) -radical over K is clearly invariant under every K-automorphism of N and hence belongs to \(N_r\) . This proves \(a\) ).

Every K-automorphism of \(\Omega\) maps N into N, hence \(N_s\) into \(N_s\) , so \(N_s\) is a quasi-Galois extension of K (V, p.~54, Prop. 3). It follows that \(N_s\) is a Galois extension of K. Every element of \(N_r \cap N_s\) is separable algebraic and \(p\) -radical over K, hence belongs to K (V, p.~Cor. 3); we thus have \(K_r \cap K_s = K\) . Now N is \(p\) -radical over \(N_s\) (V, p.~44, Prop. 13) and separable algebraic over \(N_r\) (V, p.~56, Th. 1) hence both \(p\) -radical and separable over \(K(N_r \cup N_s)\) . Hence we have \(N = K(N_r \cup N_s)\) (V, p.~39, Cor. 3) and the assertions \(b\) , \(c\) ) follow from Th. 5 (V, p.~71).

COROLLARY. --- Let p the characteristic exponent of K, \(\bar{K}\) an algebraic closure of K, \(K_{s}\) the relative separable closure of K in \(\bar{K}\) and \(K^{p^{-\infty}}\) the perfect closure of K. Then the canonical homomorphism of \(K^{p^{-\infty}} \otimes K_{s}\) into \(\bar{K}\) is an isomorphism.

Remark. --- Let R (resp. S) be a p-radical (resp. separable algebraic) extension of K. Then the algebra \(R \otimes_{K} S\) is a field : for R (resp. S) is isomorphic to a subextension of \(K^{p^{-\infty}}\) (resp. \(K_{s}\) ) and it suffices to apply the above Cor. and Prop. 1 of V, p.~17.

